% source: manti.01.ma, 1ère SM, contrôle court fonctions (Drive 0B6pfSJcEVng0dGtBakVVVWw4RjA). Ex.3 = fonctions_s5 Ex.4 -> non repris
\ex{Fonctions}{dGtB}{Ex. 1}{5}
Déterminer l'ensemble de définition de la fonction $f$ dans chacun des cas suivants :
\begin{enumerate*}[label=\arabic*),itemjoin={\quad}]
\item $f(x)=\dfrac{x+3}{|x+4|-3x}$
\item $f(x)=\dfrac{x^2}{x-\sqrt x+2}$
\item $f(x)=\dfrac{2x-3}{E(3x)-2}$
\item $f(x)=\sqrt{\dfrac x3-\dfrac3x}$
\end{enumerate*}
\fin
\ex{Fonctions}{dGtB}{Ex. 2}{5}
On considère la fonction $f$ définie sur $\mathbb{R}$ par $f(x)=\dfrac{x^2}{x^2-x+2}$.
\begin{enumerate}
\item Montrer que la fonction $f$ est majorée par $\dfrac87$.
\item La fonction $f$ admet-elle une valeur maximale ? Justifier.
\end{enumerate}
\fin
